Read in Gauge Theory Gravity: spacetime is the invertible region of the gravitational field $(\hb, \Omega)$ on flat spacetime, and the conformal spinor connection is the flat-space $\Cl(2,4)$ gauge field of Theorem \claimref{synthesis.gtg-tractor-gauge-field-reproduces-tractor-connection} acting on $\CS$ (Fact \claimref{synthesis.imported-gtg-dictionary}). A connection is given by its gauge field: along a vector $a$ it acts as $a\cdot\nabla + \mathcal{A}(a)$ in a local frame, $\mathcal{A}$ linear in $a$ with values in endomorphisms of the fibre (in forms, the connection form, a labelled reading).

A connection on a vector bundle whose parallel transports lie in a closed subgroup $G$ of the fibre's general linear group has, in a local frame adapted to $G$, a gauge field valued in the Lie algebra of $G$: the transport along a curve $\gamma$ from $\gamma(0)$ to $\gamma(t)$ is a curve in $G$ through the identity, differentiable in $t$ since the connection is smooth, and the gauge field evaluated on the tangent $\dot\gamma(0)$ is minus its derivative at $t = 0$, a tangent vector to $G$ at the identity; that a closed subgroup is an embedded Lie subgroup, so that this tangent vector lies in its Lie algebra, is the closed-subgroup theorem \cite[Thm.~20.12]{Lee2013}. For $\nabla$ the group is $\mathrm{Spin}(2,4)$, with Lie algebra $\mathfrak{so}(2,4)$, the bivectors of $\Cl(2,4)$ acting on $\CS$ by left multiplication with the factor $\tfrac12$ (Fact \claimref{standard-model.imported-tractor-bundle-spinor-bundle-their}); this half of the statement is the hypothesis read through that theorem. For $D_{\mathcal{H}}$ the transports are unitary maps of $\C^{2} \otimes \C^{3}$ carrying the set of decomposable tensors onto itself, hence of the form $A \otimes B$ with $(A, B) \in U(2) \times U(3)$ by Lemma \claimref{standard-model.segre-lemma-automorphisms-preserving-decomposable-tensors} ($m = 2 \neq 3 = k$); the group of such maps is the image of $U(2) \times U(3)$, closed since $U(2) \times U(3)$ is compact, with Lie algebra the image of $\uu(2) \oplus \uu(3)$ under $(X, X') \mapsto X \otimes 1 + 1 \otimes X'$, which is $\su(2) \oplus \su(3) \oplus \uu(1)$ acting as in Theorem \claimref{standard-model.gauge-algebra-minimal-massive-composite}. Conversely a gauge field valued in a Lie algebra $\mathfrak{g}$ has transports in the connected group generated by $\exp\mathfrak{g}$, which preserves whatever $\mathfrak{g}$ infinitesimally preserves: the spin structure for $\mathfrak{so}(2,4)$, and the decomposable tensors and the Hermitian structure for the image of $\uu(2) \oplus \uu(3)$, since $\exp(X \otimes 1 + 1 \otimes X') = e^{X} \otimes e^{X'}$.
